\section*{अध्याय \ref{ch:g11:plate-motions} --- प्लेटें और उनकी गतियाँ}

\begin{solution}{exo:g11:plate-motions:1}
\omterm{def:g11:plate-motions:rigid-plate}{दृढ़ प्लेट} \omterm{def:g11:lithosphere-asthenosphere:lithosphere}{स्थलमंडल} का ऐसा टुकड़ा है जो एक ही खंड के रूप में चलता है: उसके बिंदुओं के बीच की
दूरियाँ नहीं बदलतीं, और उसका विरूपण उसकी सीमाओं के अनुदिश
संकरे क्षेत्रों में इकट्ठा रहता है। भूकंप यह दिखाते हैं: उनमें से लगभग सभी
संकरी पट्टियों के अनुदिश, यानी प्लेट-सीमाओं पर, होते हैं, और बहुत कम
प्लेटों के भीतर।
\end{solution}

\begin{solution}{exo:g11:plate-motions:2}
प्रति वर्ष \qty{72}{mm} का अर्थ प्रति मिलियन वर्ष \qty{72}{km} है। 30 वर्षों में स्टेशन
$72 \times 30 = \qty{2160}{mm}$ चलता है, लगभग \qty{2.2}{m}; एक मिलियन वर्षों में
\qty{72}{km}।
\end{solution}

\begin{solution}{exo:g11:plate-motions:3}
$v_E = 30/10 = \qty{3}{cm}$ प्रति वर्ष और $v_N = 40/10 = \qty{4}{cm}$ प्रति वर्ष।
चाल $\sqrt{3^2 + 4^2} = \qty{5}{cm}$ प्रति वर्ष है। उत्तर से कोण $\beta$
$\tan\beta = 3/4$ को संतुष्ट करता है, इसलिए $\beta \approx \ang{37}$: स्टेशन
उत्तर-पूर्व की ओर, उत्तर से \ang{37} पूर्व, चलता है।
\end{solution}

\begin{solution}{exo:g11:plate-motions:4}
\omterm{def:g11:plate-motions:hotspot-track}{तप्त स्थल पथ} ज्वालामुखियों और समुद्री पर्वतों की वह पंक्ति है जिसे कोई प्लेट
किसी तप्त स्थल के ऊपर से चलते हुए उससे दूर ले जाती है, जो सक्रिय ज्वालामुखी से दूरी के साथ
अधिकाधिक पुराने होते हैं। सबसे नया ज्वालामुखी, किलाउआ, श्रृंखला के
दक्षिण-पूर्वी सिरे पर, हवाई द्वीप पर, है; आयु
पश्चिम-उत्तर-पश्चिम की ओर बढ़ती है। प्लेट अपने ज्वालामुखियों को तप्त स्थल से
अपनी गति की दिशा में दूर ले जाती है: प्रशांत प्लेट
पश्चिम-उत्तर-पश्चिम की ओर चलती है।
\end{solution}

\begin{solution}{exo:g11:plate-motions:5}
$519/5.1 \approx \qty{102}{km}$ प्रति मिलियन वर्ष, यानी प्रति वर्ष \qty{102}{mm},
लगभग प्रति वर्ष \qty{10}{cm}।
\end{solution}

\begin{solution}{exo:g11:plate-motions:6}
$v_E = 63/31 \approx \qty{20.3}{mm}$ प्रति वर्ष और $v_N = 48/31 \approx
\qty{15.5}{mm}$ प्रति वर्ष। चाल $\sqrt{20.3^2 + 15.5^2} \approx
\qty{25.5}{mm}$ प्रति वर्ष है; $\tan\alpha = 15.5/20.3 \approx 0.76$ से $\alpha
\approx \ang{37}$ मिलता है: स्टेशन पूर्व से \ang{37} उत्तर, यानी
उत्तर-पूर्व की ओर, चलता है। मौना की $72/25.5 \approx 2.8$ गुना तेज़ चलता है।
\end{solution}

\begin{solution}{exo:g11:plate-motions:7}
(1) मिलीमीटर प्रति वर्ष में,
\[
  \vec v_{\mathrm{NA/EU}} = (-10.5 - 7.7,\ 20.6 - 16.3) = (-18.2, +4.3).
\] इसकी चाल $\sqrt{18.2^2 + 4.3^2} \approx
\qty{18.7}{mm}$ प्रति वर्ष है, और $\tan\alpha = 4.3/18.2 \approx 0.24$ से
पश्चिम से \ang{13} उत्तर मिलता है। (2) दोनों प्लेटें कटक के आर-पार
प्रति वर्ष लगभग \qty{1.9}{cm} की दर से एक-दूसरे से दूर होती हैं: यह आइसलैंड में कटक की \omterm{def:g11:drift-to-spreading:spreading-rate}{पूर्ण प्रसरण दर}
है। (3) यह उत्तरी अटलांटिक के लिए उद्धृत प्रति वर्ष \qty{2}{cm} से
मेल खाती है।
\end{solution}

\begin{solution}{exo:g11:plate-motions:8}
(1) $v = \pi \times 0.651 \times 6371/180 \approx \qty{72.4}{km}$ प्रति मिलियन
वर्ष, लगभग प्रति वर्ष \qty{7.2}{cm}। (2) $72.4 \times \sin\ang{30} \approx
\qty{36.2}{km}$ प्रति मिलियन वर्ष, लगभग प्रति वर्ष \qty{3.6}{cm}। (3) घूर्णन
ध्रुव पर, यदि वह प्लेट पर हो। चाल बिंदु-दर-बिंदु बदलती है,
ध्रुव पर शून्य से उससे \ang{90} दूर अपने सबसे बड़े मान तक: किसी
पूरी प्लेट की कोई एक चाल नहीं होती, केवल हर बिंदु पर एक चाल (या उसकी सतह पर
एक औसत)।
\end{solution}

\begin{solution}{exo:g11:plate-motions:9}
(1) मिलीमीटर प्रति वर्ष में,
\[
  \vec v_{\mathrm{IN/EU}} = (40.6 - 24.0,\ 38.6 + 3.3) = (+16.6, +41.9).
\] चाल $\sqrt{16.6^2 + 41.9^2} \approx
\qty{45}{mm}$ प्रति वर्ष है; उत्तर से कोण $\beta$ $\tan\beta =
16.6/41.9 \approx 0.40$ को संतुष्ट करता है, इसलिए $\beta \approx \ang{22}$: भारत
यूरेशिया की ओर, उत्तर से \ang{22} पूर्व, चलता है। (2) प्रति मिलियन वर्ष \qty{45}{km}, इसलिए
10 मिलियन वर्षों में \qty{450}{km}।
\end{solution}

\begin{solution}{exo:g11:plate-motions:10}
$(2432 - 1058)/(27.7 - 10.3) = 1374/17.4 \approx \qty{79}{km}$ प्रति मिलियन
वर्ष, लगभग प्रति वर्ष \qty{7.9}{cm}, प्रति वर्ष \qty{9.1}{cm} से थोड़ा कम।
हर आयु लगभग आधे मिलियन वर्ष तक अनिश्चित है; आयु का
अंतर, $17.4 \pm 1.0$ मिलियन वर्ष, लगभग 75 और
\qty{84}{km} प्रति मिलियन वर्ष के बीच की चाल देता है। इन दोनों ज्वालामुखियों के बीच प्लेट
औसत से थोड़ा धीरे चली लगती है।
\end{solution}

\begin{solution}{exo:g11:plate-motions:11}
(1) $\vec v_{A/C} = \vec v_{A/B} + \vec v_{B/C}$: प्रति वर्ष \qty{3}{cm} पूर्व और
प्रति वर्ष \qty{4}{cm} उत्तर। दोनों तीर समकोण बनाते हैं, और त्रिभुज की तीसरी
भुजा की लंबाई $\sqrt{3^2 + 4^2} = \qty{5}{cm}$ प्रति वर्ष है,
उत्तर से \ang{37} पूर्व की दिशा में ($\tan\beta = 3/4$)। (2) $\vec v_{C/A}
= -\vec v_{A/C}$: विपरीत दिशा में प्रति वर्ष \qty{5}{cm}, दक्षिण से \ang{37}
पश्चिम।
\end{solution}

\begin{solution}{exo:g11:plate-motions:12}
(1) $54/0.375 = \qty{144}{km}$ प्रति मिलियन वर्ष, प्रति वर्ष \qty{14.4}{cm}। (2)
$29/0.375 \approx 77$ से $79/0.375 \approx \qty{211}{km}$ प्रति मिलियन
वर्ष तक: लगभग 8 और \qty{21}{cm} प्रति वर्ष के बीच। (3) तप्त स्थल के पास
दूरी किसी ज्वालामुखी के आकार से या उस क्षेत्र के आकार से बड़ी नहीं है जिस पर
तप्त स्थल ज्वालामुखियों को पोषित करता है, और आयु उस समय से बड़ी नहीं है जो किसी ज्वालामुखी को
बढ़ने में लगता है: सापेक्ष रूप से दोनों बहुत अनिश्चित हैं। तप्त स्थल से दूर
\qty{2400}{km} पर \qty{25}{km} की, या 28 मिलियन पर आधे मिलियन वर्ष की, त्रुटि
केवल एक या दो प्रतिशत है।
\end{solution}

\begin{solution}{exo:g11:plate-motions:13}
(1) बिखराव का दस गुना \qty{50}{mm} है: $50/72 \approx 0.7$ वर्ष, लगभग
आठ महीने। (2) $50/5 = 10$ वर्ष। (3) मिलियनों वर्षों में प्लेटें
दसियों से सैकड़ों किलोमीटर चलती हैं, ऐसी दूरियाँ जिन्हें चुंबकीय धारियाँ अंकित करती हैं और
जिन्हें समुद्र पर मापना आसान था; कुछ मिलीमीटर तक परिशुद्ध स्थिति-निर्धारण,
जो गति को महीनों में दिखा देता है, बहुत बाद में आया।
\end{solution}

\begin{solution}{exo:g11:plate-motions:14}
(1) मिलीमीटर प्रति वर्ष में,
\[
  \vec v_{\mathrm{SA/AF}} = (-4.9 - 27.1,\ 11.0 - 17.9) = (-32.0, -6.9).
\] चाल $\sqrt{32.0^2 + 6.9^2} \approx
\qty{32.7}{mm}$ प्रति वर्ष है, और $\tan\alpha = 6.9/32.0 \approx 0.22$ से
पश्चिम से \ang{12} दक्षिण मिलता है। (2) \omterm{def:g11:drift-to-spreading:spreading-rate}{पूर्ण प्रसरण दर} लगभग प्रति वर्ष \qty{3.3}{cm}
है, अर्ध दर लगभग प्रति वर्ष \qty{1.6}{cm}। (3) पिछले
कुछ मिलियन वर्षों की अर्ध दर, प्रति वर्ष \qty{1.6}{cm}, 66 मिलियन वर्षों पर औसत
प्रति वर्ष \qty{1.9}{cm} से थोड़ी कम है: हो सकता है दक्षिण अटलांटिक
हाल में औसत से थोड़ा धीरे खुला हो, या अंतर
अवसाद-आयुओं की अनिश्चितताओं से आया हो। दोनों मापन सहमत हैं
कि कटक धीमा है।
\end{solution}

\begin{solution}{exo:g11:plate-motions:15}
(1) प्रति वर्ष \qty{20}{mm} बढ़ती एक सरल रेखा, पाँच वर्ष बाद
\qty{50}{mm} की एकाएक गिरावट, फिर उसी ढाल से दोबारा बढ़ती रेखा।
(2) स्टेशन अंत में अपने आरंभ से $10 \times 20 - 50 = \qty{150}{mm}$ पूर्व में है:
$150/10 = \qty{15}{mm}$ प्रति वर्ष, एक-चौथाई कम। (3) हर आधा भाग
प्रति वर्ष \qty{20}{mm} देता है। (4) उछाल भूकंप से हुई एक अलग, एकाएक
गति है: भूगणितज्ञ उसका परिमाण उसके ठीक पहले और
ठीक बाद की स्थितियों से मापता है, उसे श्रेणी से हटा देता है (या दोनों भागों को
एक ही ढाल और उनके बीच एक सीढ़ी के साथ मिलाता है), और प्रति वर्ष \qty{20}{mm} को
स्टेशन के वेग के रूप में रखता है।
\end{solution}

\begin{solution}{pb:g11:plate-motions:1}
\textbf{1.} मौना की (\num{0.375} मिलियन वर्ष) सबसे नया है, सुइको
(केंद्र) (\num{64.7} मिलियन वर्ष) सबसे पुराना। आयु
किलाउआ से दूरी के साथ बढ़ती है।

\textbf{2.} $\qty{1}{km} = 10^6~\unit{mm}$ और एक मिलियन वर्ष $10^6$
वर्ष हैं: प्रति मिलियन वर्ष \qty{1}{km} का अर्थ $10^6~\unit{mm}/10^6$ वर्ष है,
यानी प्रति वर्ष \qty{1}{mm}।

\textbf{3.} $519/5.1 \approx \qty{102}{mm}$ प्रति वर्ष।

\textbf{4.} नेकर $1058/10.3 \approx 102.7$; लेसन $1818/19.9 \approx
91.4$; मिडवे $2432/27.7 \approx 87.8$; कोलाहान $3128/38.6 \approx 81.0$;
ऐबट $3280/38.7 \approx 84.8$ (मिलीमीटर प्रति वर्ष)।

\textbf{5.} बिंदु (54; 0.375), (519; 5.1), (780; 7.2), (1058; 10.3),
(1818; 19.9), (2432; 27.7), (3128; 38.6), (3280; 38.7) मूल-बिंदु से होकर जाने वाली
एक सरल रेखा के निकट हैं, कुछ बिखराव के साथ (प्रति मिलियन वर्ष \qty{91}{km} की
रेखा की तुलना में काउआई, निहोआ और नेकर उससे थोड़ा नीचे हैं, लेसन उस पर,
कोलाहान और ऐबट थोड़ा ऊपर)।

\textbf{6.} हर ज्वालामुखी तप्त स्थल के ऊपर बना और तब से
प्लेट के साथ चला है। समय $t$ में प्लेट एक नियत
दिशा में $v\,t$ चली है, इसलिए ज्वालामुखी उस बिंदु से $d = v\,t$ पर है जो अब
तप्त स्थल के ऊपर है: $t = d/v$।

\textbf{7.} $(102.7 + 91.4 + 87.8 + 81.0 + 84.8)/5 \approx \qty{89.5}{mm}$ प्रति
वर्ष।

\textbf{8.} 89.5, प्रति मिलियन वर्ष \qty{91}{km} की तुलना में: दोनों में
\qty{2}{\percent} से कम का अंतर है।

\textbf{9.} $519/91 \approx 5.7$ मिलियन वर्ष, जबकि आयु \num{5.1}
मिलियन वर्ष है: निकट, काउआई औसत रेखा के पूर्वानुमान से थोड़ा
नया है।

\textbf{10.} ज्वालामुखी पश्चिम-उत्तर-पश्चिम की ओर अधिकाधिक पुराने हैं:
प्लेट उन्हें उसी दिशा में ले जाती है। प्रशांत प्लेट तप्त स्थल के ऊपर से
पश्चिम-उत्तर-पश्चिम की ओर चलती है।

\textbf{11.} $(4860 - 3493)/(64.7 - 42.4) = 1367/22.3 \approx \qty{61}{km}$
प्रति मिलियन वर्ष, लगभग प्रति वर्ष \qty{6.1}{cm}।

\textbf{12.} $65/91 \approx 0.71$: सम्राट श्रृंखला के अनुदिश प्लेट
हवाई श्रृंखला के अनुदिश की अपेक्षा लगभग \qty{30}{\percent} धीरे चली।

\textbf{13.} $3668/50 \approx \qty{73}{km}$ प्रति मिलियन वर्ष, लगभग
प्रति वर्ष \qty{7.3}{cm}, प्रति मिलियन वर्ष \qty{91}{km} से कम। मूल-बिंदु से होकर रेखा
1987 की आयुओं का उपयोग करती है, जो हवाई श्रृंखला का अंत 39 से 43
मिलियन वर्ष पर रखती हैं; मोड़ के आरंभ की लगभग 50 मिलियन वर्ष की नई आयु
उसी दूरी को अधिक लंबे समय में फैला देती है। दोनों सही हैं यदि प्लेट
हवाई श्रृंखला के आरंभिक भाग में अधिक धीरे और बाद में अधिक तेज़ चली:
चाल नियत नहीं थी।

\textbf{14.} प्रशांत प्लेट की गति की दिशा में परिवर्तन, लगभग
50 मिलियन वर्ष पहले।

\textbf{15.} श्रृंखला की दिशा वह दिशा है जिसमें प्लेट
चली; उसके अनुदिश आयुओं का अंतराल उसकी चाल है। कोई प्लेट बिना मुड़े तेज़ या
धीमी हो सकती है, या चाल बदले बिना मुड़ सकती है, इसलिए दोनों
स्वतंत्र हैं। सम्राट
श्रृंखला के अनुदिश दिशा बदले बिना चाल के परिवर्तन दिखाते हैं कि उसका चाल-अभिलेख और मोड़
अलग-अलग घटनाएँ हैं: मोड़ तप्त स्थल के धीमे पड़ने से नहीं बना।

\textbf{16.} $\sqrt{63.2^2 + 34.5^2} \approx \qty{72.0}{mm}$ प्रति वर्ष, और
$\tan\alpha = 34.5/63.2 \approx 0.55$ से पश्चिम से \ang{29} उत्तर मिलता है।

\textbf{17.} पश्चिम-उत्तर-पश्चिम, पश्चिम से \ang{22.5} उत्तर है, उत्तर-पश्चिम
\ang{45}: स्टेशन की दिशा, पश्चिम से \ang{29} उत्तर, इन दोनों के बीच है,
श्रृंखला की दिशा के निकट। यह संगत है।

\textbf{18.} स्टेशन का वेग ऐसे निर्देश-तंत्र के सापेक्ष मापा जाता है जिसमें
\omterm{def:g11:lithosphere-asthenosphere:lithosphere}{स्थलमंडल} समग्र रूप से नहीं घूमता, जबकि श्रृंखला तप्त स्थल के सापेक्ष
गति देती है: यदि दोनों निर्देश-तंत्र एक-दूसरे के सापेक्ष चलते हैं,
तो चालें भिन्न होती हैं। और स्टेशन पिछले 30 वर्ष मापता है,
श्रृंखला पिछले 40 मिलियन वर्ष, जिनमें चाल बदली हो सकती है।
आयुओं का बिखराव भी श्रृंखला की चाल पर प्रति वर्ष कुछ मिलीमीटर की
अनिश्चितता छोड़ता है।

\textbf{19.} $100/91 \approx 1.1$ मिलियन वर्ष।

\textbf{20.} \textbf{पिछले 40 मिलियन वर्षों में प्रशांत प्लेट हवाई तप्त स्थल के ऊपर से
प्रति वर्ष लगभग \qty{9}{cm} की दर से} पश्चिम-उत्तर-पश्चिम की ओर चली है;
मोड़ से पहले वह अधिक धीरे, और दूसरी दिशा में, चलती थी।
\end{solution}
